\textbf{Phần 1. Trắc nghiệm lựa chọn}
\setcounterpageref{ex}{0}

\begin{ex} %[5D7V44]
Hai kho chứa tất cả $360$ tấn lúa. Nếu chuyển $30$ tấn từ kho A sang kho B thì số thóc kho A bằng $\dfrac{5}{7}$ số thóc kho B. Số thóc kho A ban đầu là
\choice
{$150\text{ tấn}$}
{\True $180\text{ tấn}$}
{$210\text{ tấn}$}
{$160\text{ tấn}$}
\loigiai{
Sau khi chuyển, tổng số thóc hai kho vẫn là $360$ tấn.\\
Tổng số phần bằng nhau là: $5 + 7 = 12\text{ (phần)}$.\\
Số thóc kho A sau khi chuyển là:
$360 : 12 \times 5 = 150\text{ (tấn)}$.\\
Số thóc kho A lúc đầu là:
$150 + 30 = 180\text{ (tấn)}$.\\
\textbf{Đáp án của câu hỏi là B}
}
\end{ex}

\begin{ex} %[5D7V44]
Một người mua một món hàng với giá $600\,000$ đồng rồi bán lại với giá $720\,000$ đồng. Người đó đã lãi bao nhiêu phần trăm so với tiền vốn?
\choice
{$12\%$}
{$15\%$}
{\True $20\%$}
{$25\%$}
\loigiai{
Số tiền lãi là:
$720\,000 - 600\,000 = 120\,000\text{ (đồng)}$.\\
Tỉ số phần trăm tiền lãi so với tiền vốn là:
$120\,000 : 600\,000 = 0{,}2 = 20\%$.\\
\textbf{Đáp án của câu hỏi là C}
}
\end{ex}

\begin{ex} %[5D7V44]
Trên bản đồ tỉ lệ $1 : 2\,500\,000$, khoảng cách giữa hai tỉnh đo được là $6\text{ cm}$. Một ô tô chạy từ tỉnh này sang tỉnh kia với vận tốc $50\text{ km/h}$. Thời gian ô tô chạy hết quãng đường đó là
\choice
{$2\text{ giờ}$}
{$2{,}5\text{ giờ}$}
{\True $3\text{ giờ}$}
{$3{,}5\text{ giờ}$}
\loigiai{
Khoảng cách thực tế là:
$6 \times 2\,500\,000 = 15\,000\,000\text{ (cm)} = 150\text{ km}$.\\
Thời gian ô tô di chuyển là:
$150 : 50 = 3\text{ (giờ)}$.\\
\textbf{Đáp án của câu hỏi là C}
}
\end{ex}

\begin{ex} %[5D7V44]
Một mảnh đất hình chữ nhật có chu vi $160\text{ m}$. Chiều dài hơn chiều rộng $20\text{ m}$. Tỉ số phần trăm của chiều rộng so với chiều dài là
\choice
{$40\%$}
{$50\%$}
{\True $60\%$}
{$75\%$}
\loigiai{
Nửa chu vi mảnh đất là: $160 : 2 = 80\text{ (m)}$.\\
Chiều rộng mảnh đất là: $(80 - 20) : 2 = 30\text{ (m)}$.\\
Chiều dài mảnh đất là: $30 + 20 = 50\text{ (m)}$.\\
Tỉ số phần trăm chiều rộng so với chiều dài là:
$30 : 50 = 0{,}6 = 60\%$.\\
\textbf{Đáp án của câu hỏi là C}
}
\end{ex}

\begin{ex} %[5D7V44]
Số học sinh của một trường tiểu học trong năm học vừa qua tăng từ $800$ học sinh lên $864$ học sinh. Số học sinh của trường đó đã tăng thêm bao nhiêu phần trăm?
\choice
{$6{,}4\%$}
{\True $8\%$}
{$8{,}64\%$}
{$10\%$}
\loigiai{
Số học sinh tăng thêm là:
$864 - 800 = 64\text{ (học sinh)}$.\\
Số phần trăm tăng thêm là:
$64 : 800 = 0{,}08 = 8\%$.\\
\textbf{Đáp án của câu hỏi là B}
}
\end{ex}

\begin{ex} %[5D7V44]
Hiện nay tuổi bố gấp $3$ lần tuổi con. Biết rằng $5$ năm trước bố hơn con $26$ tuổi. Tuổi của con hiện nay là
\choice
{$11\text{ tuổi}$}
{\True $13\text{ tuổi}$}
{$15\text{ tuổi}$}
{$18\text{ tuổi}$}
\loigiai{
Hiệu số tuổi của bố và con không đổi theo thời gian và luôn bằng $26$ tuổi.\\
Hiệu số phần bằng nhau hiện nay là: $3 - 1 = 2\text{ (phần)}$.\\
Tuổi con hiện nay là:
$26 : 2 = 13\text{ (tuổi)}$.\\
\textbf{Đáp án của câu hỏi là B}
}
\end{ex}

\begin{ex} %[5D7V44]
Một cửa hàng giảm giá một đôi giày $20\%$, sau đó nhân dịp khai trương lại giảm thêm $10\%$ trên giá đã giảm. So với giá niêm yết ban đầu, đôi giày đó đã được giảm tất cả bao nhiêu phần trăm?
\choice
{$30\%$}
{\True $28\%$}
{$25\%$}
{$18\%$}
\loigiai{
Sau lần giảm thứ nhất, giá đôi giày bằng: $100\% - 20\% = 80\%$ giá ban đầu.\\
Sau lần giảm thứ hai, giá đôi giày bằng: $80\% \times (100\% - 10\%) = 72\%$ giá ban đầu.\\
Tổng mức giảm giá là: $100\% - 72\% = 28\%$.\\
\textbf{Đáp án của câu hỏi là B}
}
\end{ex}

\begin{ex} %[5D7V44]
Hai số có hiệu bằng $56$. Biết số lớn gấp $5$ lần số bé. Tổng của hai số đó là
\choice
{$70$}
{\True $84$}
{$96$}
{$112$}
\loigiai{
Hiệu số phần bằng nhau là: $5 - 1 = 4\text{ (phần)}$.\\
Số bé là: $56 : 4 = 14$.\\
Số lớn là: $14 \times 5 = 70$.\\
Tổng của hai số là: $70 + 14 = 84$.\\
\textbf{Đáp án của câu hỏi là B}
}
\end{ex}

\begin{ex} %[5D7V44]
Trên bản đồ tỉ lệ $1 : 5\,000$, một khu đất hình vuông có diện tích đo được trên bản đồ là $4\text{ cm}^2$. Diện tích thực tế của khu đất đó bằng bao nhiêu héc-ta?
\choice
{$0{,}1\text{ ha}$}
{\True $1\text{ ha}$}
{$10\text{ ha}$}
{$100\text{ ha}$}
\loigiai{
Cạnh khu đất hình vuông trên bản đồ là: $2\text{ cm}$ (vì $2 \times 2 = 4$).\\
Cạnh thực tế của khu đất là:
$2 \times 5\,000 = 10\,000\text{ (cm)} = 100\text{ m}$.\\
Diện tích thực tế của khu đất là:
$100 \times 100 = 10\,000\text{ (m}^2\text{)} = 1\text{ ha}$.\\
\textbf{Đáp án của câu hỏi là B}
}
\end{ex}

\begin{ex} %[5D7V44]
Một can nước ngọt nặng $10\text{ kg}$, trong đó nước ngọt chiếm $90\%$ khối lượng, còn lại là vỏ can. Sau khi người ta dùng hết một nửa lượng nước ngọt thì cả can và lượng nước ngọt còn lại nặng bao nhiêu ki-lô-gam?
\choice
{$5\text{ kg}$}
{\True $5{,}5\text{ kg}$}
{$6\text{ kg}$}
{$4{,}5\text{ kg}$}
\loigiai{
Lượng nước ngọt ban đầu nặng: $10 \times 90\% = 9\text{ (kg)}$.\\
Khối lượng vỏ can rỗng là: $10 - 9 = 1\text{ (kg)}$.\\
Nửa lượng nước ngọt nặng: $9 : 2 = 4{,}5\text{ (kg)}$.\\
Cả can và nước ngọt còn lại cân nặng: $4{,}5 + 1 = 5{,}5\text{ (kg)}$.\\
\textbf{Đáp án của câu hỏi là B}
}
\end{ex}

\begin{ex} %[5D7V44]
Hai lớp 5A và 5B có tất cả $75$ học sinh. Nếu chuyển $3$ học sinh từ lớp 5A sang lớp 5B thì số học sinh lớp 5A bằng $\dfrac{2}{3}$ số học sinh lớp 5B. Số học sinh lớp 5A lúc đầu là
\choice
{$30\text{ học sinh}$}
{\True $33\text{ học sinh}$}
{$35\text{ học sinh}$}
{$42\text{ học sinh}$}
\loigiai{
Sau khi chuyển, tổng số học sinh của hai lớp không đổi và bằng $75$ học sinh.\\
Số học sinh lớp 5A sau khi chuyển là:
$75 : (2 + 3) \times 2 = 30\text{ (học sinh)}$.\\
Số học sinh lớp 5A lúc đầu là:
$30 + 3 = 33\text{ (học sinh)}$.\\
\textbf{Đáp án của câu hỏi là B}
}
\end{ex}

\begin{ex} %[5D7V44]
Tìm số tự nhiên $x$, biết rằng $35\%$ của số $x$ bằng $84$.
\choice
{$210$}
{\True $240$}
{$280$}
{$300$}
\loigiai{
Giá trị của $x$ là:
$84 : 35\% = 84 : 0{,}35 = 240$.\\
\textbf{Đáp án của câu hỏi là B}
}
\end{ex}

\textbf{Phần 2. Trắc nghiệm đúng sai}
\setcounterpageref{ex}{0}

\begin{ex} %[5D7V44]
Cho hai số $A$ và $B$ có tổng bằng $200$. Biết rằng $20\%$ số $A$ bằng $30\%$ số $B$.
\choiceTF
{\True Tỉ số giữa số $A$ và số $B$ là $\dfrac{3}{2}$}
{\True Tổng số phần bằng nhau là $5$ phần}
{\True Giá trị của số $A$ là $120$}
{Số $B$ lớn hơn số $A$}
\loigiai{
\begin{itemchoice}
\item Đúng. $A \times 0{,}2 = B \times 0{,}3 \implies \dfrac{A}{B} = \dfrac{0{,}3}{0{,}2} = \dfrac{3}{2}$.
\item Đúng. Tổng số phần: $3 + 2 = 5\text{ (phần)}$.
\item Đúng. Số $A$ là: $200 : 5 \times 3 = 120$.
\item Sai. Số $B$ là: $200 - 120 = 80 < 120$.
\end{itemchoice}
}
\end{ex}

\begin{ex} %[5D7V44]
Một mảnh đất hình chữ nhật có chiều dài $60\text{ m}$ và chiều rộng bằng $\dfrac{2}{3}$ chiều dài.
\choiceTF
{\True Chiều rộng của mảnh đất là $40\text{ m}$}
{\True Diện tích của mảnh đất là $2\,400\text{ m}^2$}
{\True Người ta dùng $25\%$ diện tích mảnh đất để đào ao nuôi cá, diện tích ao cá là $600\text{ m}^2$}
{Diện tích phần đất còn lại để trồng cây ăn quả là $1\,500\text{ m}^2$}
\loigiai{
\begin{itemchoice}
\item Đúng. Chiều rộng: $60 \times \dfrac{2}{3} = 40\text{ (m)}$.
\item Đúng. Diện tích: $60 \times 40 = 2\,400\text{ (m}^2\text{)}$.
\item Đúng. Diện tích ao cá: $2\,400 \times 25\% = 600\text{ (m}^2\text{)}$.
\item Sai. Diện tích trồng cây là: $2\,400 - 600 = 1\,800\text{ (m}^2\text{)} 
eq 1\,500\text{ m}^2$.
\end{itemchoice}
}
\end{ex}

\textbf{Phần 3. Trắc nghiệm trả lời ngắn}
\setcounterpageref{ex}{0}

\begin{ex} %[5D7V44]
Hai số có tổng bằng $150$. Tỉ số của hai số đó là $\dfrac{2}{3}$. Tìm số lớn.
\shortans{90}
\loigiai{
Tổng số phần bằng nhau là:
$2 + 3 = 5\text{ (phần)}$.\\
Số lớn là:
$150 : 5 \times 3 = 90$.
}
\end{ex}

\begin{ex} %[5D7V44]
Hiệu của hai số bằng $64$. Tỉ số của hai số đó là $\dfrac{3}{7}$. Tìm số bé.
\shortans{48}
\loigiai{
Hiệu số phần bằng nhau là:
$7 - 3 = 4\text{ (phần)}$.\\
Số bé là:
$64 : 4 \times 3 = 48$.
}
\end{ex}

\begin{ex} %[5D7V44]
Tìm giá trị $20\%$ của số $450$.
\shortans{90}
\loigiai{
Giá trị cần tìm là:
$450 \times 20\% = 90$.
}
\end{ex}

\begin{ex} %[5D7V44]
Một trường tiểu học có $500$ học sinh. Số học sinh tham gia câu lạc bộ cờ vua chiếm $18\%$ tổng số học sinh toàn trường. Hỏi có bao nhiêu học sinh tham gia câu lạc bộ cờ vua?
\shortans{90}
\loigiai{
Số học sinh tham gia câu lạc bộ cờ vua là:
$500 \times 18\% = 90\text{ (học sinh)}$.
}
\end{ex}

\begin{ex} %[5D7V44]
Trên bản đồ tỉ lệ $1 : 100\,000$, quãng đường từ thị xã A đến thị xã B dài $5\text{ cm}$. Khoảng cách thực tế giữa hai thị xã đó là bao nhiêu ki-lô-mét?
\shortans{5}
\loigiai{
Khoảng cách thực tế là:
$5 \times 100\,000 = 500\,000\text{ (cm)}$.\\
Đổi: $500\,000\text{ cm} = 5\text{ km}$.
}
\end{ex}

\begin{ex} %[5D7V44]
Một cửa hàng nhập về một chiếc áo với giá vốn là $150\,000$ đồng. Cửa hàng bán ra với mức lãi $20\%$ so với giá vốn. Số tiền lãi thu được từ chiếc áo đó là bao nhiêu nghìn đồng?
\shortans{30}
\loigiai{
Số tiền lãi là:
$150\,000 \times 20\% = 30\,000\text{ (đồng)}$.\\
Đổi: $30\,000\text{ đồng} = 30\text{ nghìn đồng}$.
}
\end{ex}
